% Source: https://artofproblemsolving.com/community/c3h1487997_sequences
% Stable identity: aops-topic-1487997-q01
\textbf{(a)} The first differences are $2,3,4,5,\dots$ and the second differences are all $1$. Since $u_{n+2}-2u_{n+1}+u_n=2a$, we get $a=\dfrac{1}{2}$. Then $u_2-u_1=3a+b=2$ gives $b=\dfrac{1}{2}$, and $u_1=a+b+c=1$ gives $c=0$. So
\[
u_n=\dfrac{1}{2}n^2+\dfrac{1}{2}n=\dfrac{n(n+1)}{2},
\]
the triangular numbers. Every term is an integer although $a$ and $b$ are not.

\textbf{(b)} Note the identity
\[
an^2+bn+c=2a\cdot\dfrac{n(n-1)}{2}+(a+b)\,n+c ,
\]
which holds because $2a\cdot\dfrac{n(n-1)}{2}=an^2-an$.

\emph{If $2a$, $a+b$, $c$ are integers}, then since $\dfrac{n(n-1)}{2}$ is an integer (one of $n$, $n-1$ is even), every $u_n$ is an integer.

\emph{Conversely}, suppose every $u_n$ is an integer. Then
\[
2a=u_3-2u_2+u_1,\qquad a+b=u_2-u_1-2a,\qquad c=u_1-(a+b)
\]
are integers (the first because $u_3-2u_2+u_1=(9a-8a+a)+(3b-4b+b)+(c-2c+c)=2a$, the second because $u_2-u_1=3a+b$, the third because $u_1=a+b+c$). $\blacksquare$

For the triangular numbers, $2a=1$, $a+b=1$, $c=0$.
